%! AP Statistics — Unit 9, Lesson 9.2 — Group Activity
\documentclass[10pt]{article}
\usepackage{apstats-article}
\usepackage{apstats-boxes}

\begin{document}

\pageheader{Unit 9, Lesson 9.2}{Group Activity}
\namepartnerperiod

\vspace{0.1in}
\textbf{Building a CI for the Slope --- Car Weight and Fuel Efficiency}

\vspace{0.05in}
A researcher randomly selects 30 cars from a large database and records curb weight
($x$, in hundreds of pounds) and highway fuel efficiency ($y$, mpg).
Computer output is shown below.

\vspace{0.1in}
\begin{center}
\renewcommand{\arraystretch}{1.3}
\begin{tabular}{lrrrr}
    \textbf{Predictor} & \textbf{Coef} & \textbf{SE Coef} & \textbf{T} & \textbf{P} \\
    \hline
    Constant & 47.8 & 3.4 & 14.06 & 0.000 \\
    Weight   & $-2.3$ & 0.41 & $-5.61$ & 0.000 \\
    \hline
    \multicolumn{5}{l}{$s = 3.2$ \quad $R^2 = 53.2\%$ \quad $n = 30$}
\end{tabular}
\end{center}

\vspace{0.05in}
The researcher wants to build a 95\% CI for the true slope $\beta$.
Use $t^* = 2.048$ ($df = 28$).

\vspace{0.1in}

\begin{tcolorbox}[colback=white, colframe=black!40,
    title=\textbf{Tier R --- Remediate},
    fonttitle=\bfseries, arc=2mm, left=3mm, right=3mm, top=2mm, bottom=2mm]

\begin{enumerate}[label=\alph*.]

    \item Identify the slope estimate and standard error from the output.

          $b = $ \blank{2.5cm}\quad and\quad $SE_b = $ \blank{2.5cm}

    \item Complete the confidence interval.

          $b \pm t^* \cdot SE_b
          = $ \blank{2cm} $\pm\ 2.048 \times $ \blank{2cm}
          $= $ \blank{2cm} $\pm$ \blank{2cm}
          $= ($ \blank{2.5cm},\; \blank{2.5cm}$)$

    \item Fill in the interpretation.

          ``We are 95\% confident that the true slope relating
          \blank{3.5cm} to \blank{3.5cm}
          for \blank{3.5cm}
          is between \blank{2cm} and \blank{2cm}
          mpg per hundred pounds.''

\end{enumerate}
\end{tcolorbox}

\begin{tcolorbox}[colback=white, colframe=black!40,
    title=\textbf{Tier A --- Approaching Proficiency},
    fonttitle=\bfseries, arc=2mm, left=3mm, right=3mm, top=2mm, bottom=2mm]

\begin{enumerate}[label=\alph*.]

    \item Identify the slope estimate $b$ and its standard error $SE_b$ from the output.
          \writelines{1}

    \item Verify the LINER conditions for this context.
          State how each condition is addressed.
          \writelines{5}

    \item Construct the 95\% CI for $\beta$. Show all work.
          \writelines{3}

    \item Interpret the CI in context. Write a complete sentence.
          \writelines{2}

    \item Does the CI provide evidence that $\beta \ne 0$? Justify your answer.
          \writelines{2}

\end{enumerate}
\end{tcolorbox}

\begin{tcolorbox}[colback=white, colframe=black!40,
    title=\textbf{Tier E --- Extension},
    fonttitle=\bfseries, arc=2mm, left=3mm, right=3mm, top=2mm, bottom=2mm]

Complete all of Tier A above. Then answer the following.

\begin{enumerate}[label=\alph*., start=6]

    \item If the researcher had collected data from only 15 cars instead of 30, how
          would the width of the CI change? Explain using the formula for $SE_b$.
          \writelines{3}

    \item A car manufacturer claims $\beta = -2.0$ (each 100 lb increase reduces
          efficiency by 2 mpg). Is this consistent with the 95\% CI? Explain.
          \writelines{3}

\end{enumerate}
\end{tcolorbox}

\end{document}
